
\section{谱序列的代数构造}
\label{sec:6-3}

本节的目的，
是从滤过微分 $\Lambda_{0,\mathrm{nov}}$-模
构造谱序列。
为此先准备若干代数材料。
本节用 $\Lambda_{0,\mathrm{nov}}$
简记 $\Lambda_{0,\mathrm{nov}}(R)$，
其中 $R$ 是任意域。

\subsection{c.f.z.}
\label{subsec:6-3-1}

先考虑以下情形。
设
$V=(\Lambda_{0,\mathrm{nov}}^{(0)})^{\oplus I}$
是自由 $\Lambda_{0,\mathrm{nov}}^{(0)}$-模，
不要求有限生成。
这里 $\Lambda_{0,\mathrm{nov}}^{(0)}$
是 $\Lambda_{0,\mathrm{nov}}$ 的零次部分。
在 $V$ 上定义滤过
\[
F^\lambda V
=\bigl(F^\lambda
\Lambda_{0,\mathrm{nov}}^{(0)}\bigr)^{\oplus I}.
\]
它在 $V$ 上诱导一个拓扑。
令 $\widehat V$ 为 $V$ 的完备化。
称这样的 $\widehat V$ 为
\emph{零能生成完备自由滤过
$\Lambda_{0,\mathrm{nov}}^{(0)}$-模}，
简称 \emph{c.f.z.}\liftIndex{I223}。
若 $V$ 还作为
$\Lambda_{0,\mathrm{nov}}^{(0)}$-模有限生成，
则称其为有限 c.f.z.。
定义称为\emph{能量函数}的映射
\[
E:\widehat V\setminus\{0\}
\longrightarrow\mathbb R_{\ge0},
\]
满足\translatornote{原书实体页 376
虽把能量函数定义在完备化
$\widehat V$ 上，却在归属式中
沿用 $F^\lambda V$。
对 $\widehat V$ 的元素须使用
原滤过诱导的闭滤过；
这里写作 $F^\lambda\widehat V$，
见疑误 A-0363。}
\[
\mathbf v\in F^{E(\mathbf v)}\widehat V,
\qquad
\mathbf v\notin F^\lambda\widehat V
\quad\text{若 }\lambda>E(\mathbf v).
\]
令
\[
\overline V
=V/\Lambda_{0,\mathrm{nov}}^{+,(0)}V
\cong R^I.
\]
我们总选取一个嵌入（分裂）
$\overline V\subset V$，
把它视作 $V$ 的零能部分，
使其与投影
$V\to\overline V$ 的复合是恒等映射。

设 $\mathbf v\in V$，写作
\[
\mathbf v=\sum_iT^{\lambda_i}v_i,
\]
其中 $v_i\in\overline V$，
$\lambda_i<\lambda_{i+1}$，
$\lim_{i\to\infty}\lambda_i=\infty$，
且 $v_i\ne0$。
称 $T^{\lambda_i}v_i$ 为 $\mathbf v$ 的\emph{分量}，
$T^{\lambda_1}v_1$ 为\emph{首分量}，
$v_1$ 为\emph{首系数}；
记首系数 $v_1$ 为 $\sigma(\mathbf v)$。
$\Lambda_{0,\mathrm{nov}}^{(0)}$
中元素的首分量和首系数亦如此定义。
注意，$\sigma(\mathbf v)$
与分裂 $\overline V\subset V$ 的选择无关。

\begin{definition}{6.3.1}\label{def:6-3-1}
设有限集
$\{\mathbf v_i\}_{1\le i\le m}\subset\widehat V$
生成子模 $W\subset\widehat V$。
若置 $E(\mathbf v_i)=\lambda_i$，
并满足
\begin{enumerate}[label=\textup{(6.3.1.\arabic*)},leftmargin=5.8em]
\item $\lambda_i\le\lambda_{i+1}$；
\item $\sigma(\mathbf v_1),\ldots,
\sigma(\mathbf v_m)$ 在 $R$ 上线性无关，
\end{enumerate}
则称此集为 $W$ 的\emph{标准基}。
\end{definition}

若 $W$ 有标准基，
显然它是自由
$\Lambda_{0,\mathrm{nov}}^{(0)}$-模。

\begin{lemma}{6.3.2}\label{lem:6-3-2}
$\widehat V$ 的任意有限生成子模 $W$
都有标准基。
\end{lemma}

\begin{proof}
设 $W$ 由 $m$ 个元素
$\mathbf v_i$（$i=1,\ldots,m$）生成，
并令 $E(\mathbf v_i)=\lambda_i$。
必要时重新排序，
可设 $\lambda_i\le\lambda_{i+1}$。
我们将修改这组元素，
使其成为标准基。
先取 $k$，使
\[
\lambda_1=\cdots=\lambda_k
<\lambda_{k+1}.
\]
对 $j=1,\ldots,k$ 的首系数
$\sigma(\mathbf v_j)$，
可找到一个元素属于 $R$ 的可逆矩阵
$(a_{i,i'})$，使线性组合
\[
\sum_{i'=1}^{k}a_{i,i'}
\sigma(\mathbf v_{i'})
\]
在 $i=1,\ldots,\ell\le k$ 时
线性无关，
在 $i>\ell$ 时为零。
用此矩阵变换部分发生元
$\mathbf v_1,\ldots,\mathbf v_k$，
便得到仍生成 $W$ 的一组
$\{\mathbf v_i\}$，满足 (6.3.3.1)。
这里条件 (6.3.3.$n$) 定义如下：

\noindent\textup{(6.3.3.$n$)}\quad
以
\[
\lambda_1=\cdots=\lambda_{k_1}
<\lambda_{k_1+1}=\cdots=\lambda_{k_2}
<\lambda_{k_2+1}=\cdots
\]
定义整数 $k_\nu$ 时，
$\sigma(\mathbf v_j)$
（$j=1,\ldots,k_n$）线性无关。

对 $n$ 归纳修改给定的
$\mathbf v_i$，使 (6.3.3.$n$) 成立。
设它已成立。
对 $V$ 施以滤过自同构后，
可设\translatornote{原书实体页 377
在式 (6.3.4) 后及紧随其后的
$X$ 生成元范围中，
两次把整数下标终点 $k_n$
误写为实数能量 $\lambda_{k_n}$。
此处按同页 (6.3.3.$n$) 的分组
统一为 $k_n$；见疑误 A-0362。}
\begin{equation}
\mathbf v_j=T^{\lambda_j}
\sigma(\mathbf v_j)
\qquad(j=1,\ldots,k_n).
\tag{6.3.4}\label{eq:6-3-4}
\end{equation}
令 $X\subset\overline V$
为 $\sigma(\mathbf v_j)$
（$j=1,\ldots,k_n$）
生成的 $R$-线性子空间。
分两种情形：

\noindent\textup{(6.3.5.1)}\quad
存在 $i\in\{k_n+1,\ldots,k_{n+1}\}$，
使 $\sigma(\mathbf v_i)\notin X$。

\noindent\textup{(6.3.5.2)}\quad
对所有 $i\in\{k_n+1,\ldots,k_{n+1}\}$，
都有 $\sigma(\mathbf v_i)\in X$。

先设 (6.3.5.1) 成立。
则存在可逆矩阵
$(a_{i,i'})$（$i,i'=k_n+1,\ldots,k_{n+1}$）
及 $a_{i,j},b_{i,j}\in R$
（$i=k_n+1,\ldots,k_{n+1}$，
$j=1,\ldots,k_n$），
使下列条件 (6.3.6) 成立；
令
\[
v'_{i,1}
=\sum_{j=1}^{k_n}b_{i,j}\sigma(\mathbf v_j)
+\sum_{i'=k_n+1}^{k_{n+1}}
a_{i,i'}\sigma(\mathbf v_{i'}).
\]

\noindent\textup{(6.3.6.1)}\quad
存在 $k_n<\ell\le k_{n+1}$，
使
\[
\{\sigma(\mathbf v_j)\}_{1\le j\le k_n}
\cup\{v'_{i,1}\}_{k_n+1\le i\le\ell}
\]
线性无关。

\noindent\textup{(6.3.6.2)}\quad
当 $i>\ell$ 时，$v'_{i,1}=0$。

现在对 $i=k_n+1,\ldots,k_{n+1}$，
用下式替换 $\mathbf v_i$：
\[
\mathbf v'_i
=\sum_{j=1}^{k_n}
b_{i,j}T^{\lambda_{k_{n+1}}-\lambda_j}
\mathbf v_j
+\sum_{i'=k_n+1}^{k_{n+1}}
a_{i,i'}\mathbf v_{i'}.
\]
注意，$k_n+1\le i\le\ell$ 时，
$\mathbf v'_i$ 的首系数是 $v'_{i,1}$。
由 (6.3.6.2)，
$\ell<i\le k_{n+1}$ 时
$E(\mathbf v'_i)>\lambda_{k_{n+1}}$。
再由 (6.3.6.1) 得到
(6.3.3.$n+1$)，
故第一种情形的归纳步成立。

若 (6.3.5.2) 成立，
则对每个 $k_n+1\le i\le k_{n+1}$，
存在 $b_{i,j}$，使
\[
v'_{i,1}=\sum_{j=1}^{k_n}
b_{i,j}\sigma(\mathbf v_j).
\]
对这些 $i$，
以如下新元替换 $\mathbf v_i$：
\[
\mathbf v'_i
=\mathbf v_i
-\sum_{j=1}^{k_n}
b_{i,j}T^{\lambda_{k_{n+1}}-\lambda_j}
\mathbf v_j
\]
在原能级 $\lambda_{k_{n+1}}$ 上，
$\mathbf v'_i$ 的系数于是变为
\[
\sigma(\mathbf v_i)
-\sum_{j=1}^{k_n}
b_{i,j}\sigma(\mathbf v_j)
=\sigma(\mathbf v_i)-\sigma(\mathbf v_i)=0.
\]
因此对全部
$k_n+1\le i\le k_{n+1}$，
都有 $E(\mathbf v'_i)>\lambda_{k_{n+1}}$。
按能量重新排列
$i=k_n+1,\ldots$ 的 $\mathbf v_i$，
再将情形分为 (6.3.5.1)、(6.3.5.2)，
重复此过程。
若遇到 (6.3.5.1)，
便进入归纳的下一步；
否则继续。
只有当全部
$\mathbf v_i$（$i=k_n+1,\ldots$）
的分量都属于 $X$ 时，
此过程才会无限进行
（由 (6.3.4) 可证）。
这时可把这些 $\mathbf v_i$
从发生元组中全部删去，
而不改变所生成的模。
引理 6.3.2 得证。
\end{proof}

由引理 6.3.2 及其前面的说明，
立即得到以下推论。

\begin{corollary}{6.3.7}\label{cor:6-3-7}
$\widehat V$ 的任意有限生成子模 $W$
都是自由且闭的。
\end{corollary}

引理 6.3.2 的证明还表明：

\begin{lemma}{6.3.2bis}\label{lem:6-3-2bis}
设 $W$ 如引理 6.3.2，
$\mathbf v_1,\ldots,\mathbf v_m$
是 $W$ 的一组发生元，
满足 $E(\mathbf v_i)\le E(\mathbf v_{i+1})$，
且
$\{\sigma(\mathbf v_1),\ldots,
\sigma(\mathbf v_{k_1})\}$
在 $R$ 上线性无关。
则可取标准基
$\mathbf v'_1,\ldots,\mathbf v'_{m'}$，
使 $i=1,\ldots,k_1$ 时
$\mathbf v_i=\mathbf v'_i$。
\end{lemma}
